Evaluation uses real marine radar logs from the MOANA \cite{moana2025} and PoLaRIS \cite{choi2024_polaris} datasets, featuring explicit environment labels and strict splits by route/location to avoid leakage. We employ MRISNet \cite{ma2024_mrisnet} as the primary downstream perception model to evaluate ship instance segmentation performance. We compare five settings: (1) fixed simulator defaults, (2) uniform random domain randomization, (3) Bayesian optimization tuning, (4) RL tuning without environment labels, and (5) RL tuning with environment labels.

The primary outcome is real-data task performance and corresponding reduction in sim-to-real gap. Secondary outcomes are cross-domain generalization (train on two environment types, test on the held-out type), robustness on unseen locations, run-to-run variance, and behavior under low-SNR/high-clutter stress tests. We also report sample efficiency as the number of real-validation calls needed to reach a target accuracy.

Each novelty claim is tied to a direct test: (a) benefit of sequential RL adaptation versus BayRn under matched evaluation budgets \cite{muratore2020bayrn}; (b) benefit of environment semantics via conditioned versus unconditioned RL; and (c) benefit of physical realism constraints via bound-violation and instability rates with/without realism regularization. This structure makes it explicit how the proposed method differs from existing randomization and adaptation families \cite{tobin2017domain,peng2018sim,chebotar2018simopt,ramos2019bayessim,openai2019rubik,josifovski2024cdr}.

Prior published results already motivate this direction (Fig.~\ref{fig:published-evidence}). In real FMCW transfer, simulation-only baselines can approach chance-level behavior, while calibrated adaptation improves substantially \cite{trinh2026digitaltwin}. In synthetic-to-measured SAR transfer, synthetic-only and simple-noise baselines are again weaker than stronger adaptation methods \cite{kim2024soft}. Together, these patterns motivate structured, feedback-driven simulator control rather than fixed randomization.
